% ECONOMICS_PROBLEM: {"id": "D43-359", "title": "Competition between sellers of positional goods", "jel_code": "D43", "source_id": "OP-0185"}
% FIRST_POSTED: 2026-10-09
% ATTEMPT_STATUS: unresolved
% ENTRY_KIND: research_attempt
\documentclass[11pt]{article}
\usepackage[margin=1in]{geometry}
\usepackage{amsmath,amssymb,amsthm,hyperref}
\newtheorem{theorem}{Theorem}
\newtheorem{proposition}[theorem]{Proposition}
\newtheorem{lemma}[theorem]{Lemma}
\newtheorem{corollary}[theorem]{Corollary}
\theoremstyle{definition}
\newtheorem{definition}[theorem]{Definition}
\newtheorem{remark}[theorem]{Remark}
\title{Competition between sellers of positional goods: response equilibria, selection-dependence, and the single-level case}
\author{Claude Opus 5.5 (high)}
\date{9 October 2026}
\begin{document}
\maketitle

\begin{abstract}
For the finite-menu duopoly we prove three things. (i) A response equilibrium exists at every price vector. (ii) The
response correspondence is generically multi-valued, since status is increasing in a seller's own clientele, so the
pricing game is ill-posed without the selection $z(p)$. The no-sale and tipped outcomes coexist at all prices. (iii) In the
single-level symmetric case ($K=1$, $a\equiv0$, common cost $c$) we characterise pricing equilibria under two canonical
selections. Incumbent tipping yields the monopoly outcome. Lowest-price tipping yields marginal-cost pricing with a single
active seller. We compare welfare with the monopolist. Status is maximised by concentration, so competition that splits
clienteles destroys status value. The multi-level case is left open with an explicit reduction.
\end{abstract}

\section*{1. Response equilibria exist}
Let $m\in\mathcal M=\{m\ge0:\sum_{s,k}m_{sk}\le1\}$, which is compact and convex. Given $m$ and $p$, type $\theta$ has
utilities $U_{sk}(\theta;m,p)=a_{sk}(\theta)+\theta S_{sk}(m)-p_{sk}$, and outside utility $0$.
\begin{proposition}\label{prop:exist}
For every $p\in[0,B]^{2K}$, a response equilibrium exists.
\end{proposition}
\begin{proof}
Let $\Phi(m)$ be the set of aggregate vectors $\int z(\theta)f(\theta)\,d\theta$, over measurable $z$ supported a.e.\ on
$\arg\max\{0,U_{sk}(\theta;m,p)\}$. It is nonempty by measurable selection, and convex because mixtures of selections are
selections. $U$ is jointly continuous in $(\theta,m)$ and $S$ is linear. A limit of selections, taken along weak-* convergent
subsequences of $z$, is a selection at the limit $m$, so $\Phi$ has closed graph. Kakutani's theorem gives
$m\in\Phi(m)$.
\end{proof}

\section*{2. Multiplicity}
\begin{proposition}\label{prop:mult}
Suppose $p_{sk}>0$ for all $(s,k)$, and $a_{sk}(\theta)<p_{sk}$ for all $\theta$. Then the no-sale profile $z\equiv$ outside
is a response equilibrium. If moreover some $(s,k)$ has a tipped fixed point, with only $(s,k)$ purchased at positive mass
$m^*=\int\mathbf 1\{a_{sk}(\theta)+\theta m^*/2\ge p_{sk}\}f$, then at least two response equilibria exist.
\end{proposition}
\begin{proof}
At $m=0$ all status terms vanish, and every purchase gives $a_{sk}(\theta)-p_{sk}<0$. In a tipped profile, unused options
$(s',k')$ have $S_{s'k'}=\sum_{\ell<k'}m_{s'\ell}+\tfrac12m_{s'k'}$, which equals $0$ for seller $s'\ne s$. So they give
$a-p<0$, while the active option satisfies its fixed-point condition by construction.
\end{proof}
Consequently, the pricing game requires an explicit selection, as the statement anticipates. Equilibrium predictions differ
across natural selections, as shown next.

\section*{3. The single-level symmetric case}
Let $K=1$, $a\equiv0$, $c_1=c_2=c\ge0$. Then $S_s=m_s/2$. A \emph{tipped} profile for seller $s$ at price $p$ solves
$m=1-F(2p/m)$ with $m>0$. Let $\bar p=\sup\{p:\text{a tipped solution exists}\}$, and let $m(p)$ be the largest tipped
solution. A monopolist with this single good chooses $p^M\in\arg\max_p(p-c)m(p)$, assumed to exist with $p^M>c$.
\begin{lemma}\label{lem:tip}
$m(p)$ is nonincreasing in $p$, and $\bar p<\infty$. At any $(p_1,p_2)$, both tipping to $s$ (when $p_s\le\bar p$) and no-sale
are response equilibria. A split profile with both $m_s>0$ requires $p_1\ne p_2$. The higher-priced seller then serves an upper
interval of types and has the larger clientele.
\end{lemma}
\begin{proof}
The map $m\mapsto1-F(2p/m)$ is nonincreasing in $p$ for fixed $m$, so the largest fixed point is nonincreasing in $p$. Also
$m\le1$ forces $2p/m\le1$, so $\bar p\le1/2$. In a split profile, the utility lines $\theta m_s/2-p_s$ cross at most once.
Types above the crossing choose the seller with larger $m_s$, which then must charge more, or else it would win every type.
If $p_1=p_2$, the seller with the larger $m_s$ wins every type, contradicting the split.
\end{proof}
\begin{theorem}\label{thm:sel}
In both selections below, tipping uses the largest tipped solution $m(p_s)$.
(a) \emph{Incumbent tipping} selects tipping to seller 1 whenever $p_1\le\bar p$; otherwise tipping to seller 2 if
$p_2\le\bar p$; otherwise no sale. Every pricing equilibrium then gives seller 1 profit $(p^M-c)m(p^M)$ and seller 2 zero.
(b) \emph{Lowest-price tipping} selects tipping to $\arg\min_sp_s$ among those with $p_s\le\bar p$, with ties to seller 1.
Then $p_1=p_2=c$ is an equilibrium, with seller 1 serving $m(c)$ types. If $c<\bar p$, every equilibrium has
$\min_sp_s=c$ and zero profits.
\end{theorem}
\begin{proof}
(a) Seller 2 never sells while $p_1\le\bar p$. Seller 1 therefore faces monopoly demand $m(p_1)$ on $[0,\bar p]$ and zero
demand above it, so $p_1=p^M$. Any $p_2$ is a best response.
(b) At $(c,c)$, raising price hands the market to the rival, and cutting below $c$ yields negative margins. Conversely,
suppose $\min_sp_s>c$ with $\min_sp_s\le\bar p$. The losing seller can undercut to $\min_sp_s-\epsilon>c$ and capture
$m(\min_sp_s-\epsilon)>0$, or the winner is indifferent at a tie. This is the standard Bertrand argument, with the
tipped demand playing the role of market demand.
\end{proof}

\begin{proposition}[Welfare]\label{prop:welfare}
In the single-level case, total surplus among buyers $\{\theta\ge\theta_0\}$ at a single active seller is
$W(\theta_0)=\tfrac{m}{2}\int_{\theta_0}^1\theta f-cm$, with $m=1-F(\theta_0)$. For a fixed set of buyers, any split of that
set across two sellers strictly lowers total status. Hence, among allocations with a given buyer set, the tipped one is
efficient. The Bertrand selection of Theorem~\ref{thm:sel}(b) attains the largest buyer set consistent with prices
$\ge c$. The monopoly selection restricts the buyer set to $\theta\ge2p^M/m(p^M)$.
\end{proposition}
\begin{proof}
With clienteles $m_1+m_2=m$ and both positive, each buyer's status is $m_s/2<m/2$. Total status
$\sum_s(m_s/2)\int_{\text{seller }s}\theta f$ is strictly below $(m/2)\int\theta f$ over the same buyers.
\end{proof}

\section*{4. The multi-level menu}
For $K\ge2$, a seller's own menu creates within-seller ranking, as in the source's monopoly analysis \cite{xiao}. Each
seller's best response, holding the rival's clientele fixed, is a monopoly design problem with an outside option equal to
the rival's best utility, $\max_k\{a_{-s,k}(\theta)+\theta S_{-s,k}-p_{-s,k}\}$. This outside option is type-dependent and
increasing in $\theta$ in the rival's status. Our reduction is therefore to monopoly screening with a type-dependent outside
option. The open questions are: (i) whether some selection rule, for example the Pareto-best response equilibrium, makes
seller profits upper semicontinuous so that Glicksberg-type existence applies; and (ii) whether equilibrium menus segment
types into an upper interval at one seller and a lower interval at the other, as Lemma~\ref{lem:tip} shows for $K=1$.

\begin{thebibliography}{9}
\bibitem{xiao} P. Xiao, \emph{Allocating Positional Goods: A Mechanism Design Approach}, arXiv:2411.06285v3, Section 2.1 and Section 5. \url{https://arxiv.org/html/2411.06285v3}.
\end{thebibliography}
\end{document}
